\documentclass[journal,onecolumn,11pt]{IEEEtran}

\usepackage{amsmath,amssymb,amsfonts}
\usepackage{algorithmic}
\usepackage{graphicx}
\usepackage{textcomp}
\usepackage{xcolor}
\usepackage{booktabs}
\usepackage{multirow}
\usepackage{siunitx}
\usepackage{hyperref}
\usepackage{cleveref}

\hypersetup{
  colorlinks=true,
  linkcolor=blue,
  citecolor=blue,
  urlcolor=blue
}

% All paper figures live under figures/. The fallback paths are
% kept so the source can be rebuilt against the original sources
% when an upstream figure is regenerated; remove them if you only
% want figures/ to be authoritative.
\graphicspath{%
  {figures/}%
  {../paper/figures/}%
  {../dissertation/figures/}%
  {../../results/r70_paper_figures/}%
}

\begin{document}

\title{Multiplicative-Gating Eleven-Axis Evaluation and Recurrent
Policy Discovery for Multi-Agent Virtual-Synchronous-Generator
Inertia/Damping Control on a Phasor-Equilibrium Backend}

\author{Yuang~Wei%
\thanks{Manuscript prepared 2026-05-18. This work is a follow-up to a
prior reproduction study \cite{wei2026reproduction} on the same
benchmark; full claim ledger, round logs, and audit reports are
released under \texttt{memory/} and \texttt{artifacts/paper\_r77/}.}}

\markboth{IEEE Transactions on Power Systems / Smart Grid, Submission}%
{Wei: Gated Eleven-Axis Evaluation and Recurrent Policies for VSG Control}

\maketitle

% ─────────────────────────────────────────────────────────────────────
\begin{abstract}
Headline geometric means hide single-axis failure modes that are
visually obvious in the published time-domain figures of a
multi-agent virtual-synchronous-generator (VSG) inertia/damping
controller. We report three contributions on the modified Kundur
four-machine benchmark of Yang \emph{et al.} (IEEE TPWRS 2023) using
the ANDES phasor-equilibrium simulator. First, we introduce a
\emph{multiplicative-gating eleven-axis ranker} (v3.1) that augments
the eight continuous frequency-and-action axes of a paper-faithful
score with three gating axes -- agent activity, late-window
oscillation, and per-agent power-balance -- whose geometric mean is
multiplied into the headline. This collapses the gap between the
analytic ranking and the paper-figure-equivalent visual ranking
(canonical-best controller v3.1 = 0.391 with $P_{\text{balance}} =
0.96$, vs.\ a pre-fix top-of-leaderboard agent that visibly collapsed
in the time-domain plot). Second, we discover a recurrent multi-agent
policy: \emph{TD3+LSTM} with an episode-keyed sequence replay and a
twenty-episode learning-rate warm-up reaches a new single-seed peak
of v3.1 = 0.430 on seed 59 ($+4.9\%$ over the previous tau-only
SOTA), with a six-seed mean of 0.369 on the healthy seed set
$\{50, 52, 54, 55, 56, 59\}$ (the broader healthy set
$\{50, 51, 52, 54, 55, 56, 59\}$ excludes only the warm-up-20
collapse of seed 51, which we treat as a per-seed warm-up
sensitivity rather than a drift failure). Third, we report a clean
\emph{negative} result for inference-time ensembling: a
heterogeneous-actor weighted ensemble (HAWE) of the top-2, top-3,
and full 6-seed actor pool consistently underperforms the single-best
agent by 2.1\% to 8.0\%; we attribute the regression to
recurrent-hidden-state divergence and contrast the negative with the
$+2$--$5\%$ HAWE gains we observed for memoryless TD3/SAC. We
additionally disclose a 44\% seed-drift rate across twelve random
seeds tested at warm-up~20 and recommend an explicit healthy-seed
filter for all downstream comparisons. Code, traces, and the eleven
referenced claim records are released under
\texttt{memory/rounds/R\{67--77\}/}.
\end{abstract}

\begin{IEEEkeywords}
Multi-agent reinforcement learning, virtual synchronous generator,
recurrent policy, LSTM, ensemble averaging, evaluation methodology,
Goodhart's law, multiplicative gating, paper-figure equivalence,
ANDES, Kundur benchmark.
\end{IEEEkeywords}

% ─────────────────────────────────────────────────────────────────────
\section{Introduction}\label{sec:intro}

\IEEEPARstart{V}{irtual} synchronous generators (VSGs)
\cite{zhong2011vsg,beck2007vsg,darco2015vsm} emulate
synchronous-machine inertia $H$ and damping $D$ at converter-interfaced
resources to slow the rate-of-change-of-frequency in low-inertia
power systems. Yang \emph{et al.}~\cite{yang2023ddic} proposed a
distributed deep-inertia controller (DDIC): a four-agent soft
actor-critic (SAC) \cite{haarnoja2018sac} framework in which each VSG
adjusts $\Delta H \in [-100, +300]$ and $\Delta D \in [-200, +600]$
based on its local frequency and that of two communication neighbours.
On the modified Kundur four-machine benchmark
\cite{kundur1994power,klein1991study} their reported DDIC/adaptive
cumulative-frequency-reward ratio is $0.62$, supporting the
multi-agent framing as a structural contribution.

A reproduction study on the ANDES phasor-equilibrium
simulator~\cite{cui2021andes,milano2018foundations} found the
single-headline geometric mean of paper-faithful axes (six axes in
the reproduction, including paper Figs.~6--9 settling-time, peak
$|\Delta f|$, action utilisation, and reward-improvement against the
no-control baseline) compresses too much information into one
scalar~\cite{wei2026reproduction}. Specifically, the highest-ranked
controller under the six-axis geometric mean turned out, on
inspection of its time-domain trace, to have an agent that
\emph{never actuated} (ES4 produced essentially zero $\Delta P$ for
the entire load-step response) while three other agents fully
absorbed the disturbance. The visual paper-figure equivalence of such
an agent is poor, yet the headline numbered placed it at rank~1.

This paper makes three contributions that together address the
single-headline gaming problem and exploit the policy-class freedom
the original work did not test.

\begin{enumerate}
  \item \textbf{Multiplicative-gating eleven-axis ranker} (v3.1,
        Sec.~\ref{sec:ranker}). We add three gating axes --
        \texttt{agent\_min\_activity}, \texttt{late\_oscillation\_inv},
        and \texttt{agent\_P\_balance} -- whose geometric mean is
        \emph{multiplied} into the eight-axis continuous geometric
        mean. A controller can only earn a high headline if every
        agent participates, oscillation has decayed by the late
        window, and per-agent power injection is balanced. We
        demonstrate the resulting ranking matches paper-figure
        visual quality (the previously-rank-1 collapsing controller
        falls below threshold).
  \item \textbf{Recurrent multi-agent policy: TD3+LSTM}
        (Sec.~\ref{sec:lstm}). We replace the static feed-forward
        actor with a TD3~\cite{fujimoto2018td3} agent whose
        actor and twin critics carry an
        \texttt{LSTMCell}~\cite{hausknecht2015drqn} hidden state
        across the episode, gated by an episode-keyed sequence
        replay buffer in the form of~\cite{ni2022recurrent}. With a twenty-episode learning-rate
        warm-up and $\tau = 0.001$, single-seed performance reaches
        v3.1 = 0.430 (CLM-0131; +4.9\% over the previous tau-only
        SOTA), with a six-seed mean of 0.369 on the healthy seed
        subset. We additionally observe that one seed (s51) peaks at
        warm-up 10 rather than warm-up 20; the per-seed warm-up
        sensitivity is reported in Table~\ref{tab:warmup}.
  \item \textbf{Negative result for inference-time HAWE on recurrent
        policies} (Sec.~\ref{sec:hawe}). A heterogeneous-actor
        weighted ensemble (HAWE) of the top-2 and top-3 healthy
        actors regresses the headline by 2.1\% to 8.0\% against the
        single best agent (CLM-0132). The negative cleanly contrasts
        with the $+2$--$5\%$ HAWE gains we observed in earlier work
        on memoryless TD3 and SAC actors, and we attribute it to
        hidden-state divergence across the ensembled actors
        (Sec.~\ref{sec:hidden-state}).
\end{enumerate}

In support of these contributions we additionally document a
\emph{seed-drift discovery}: of twelve random training seeds tested at
warm-up 20, five collapsed to a degenerate policy with LS1
$\Delta f$-tracking score below 0.05 (CLM-0133). We argue that the
$\sim$44\% drift rate is part of the controller's reported quality and
recommend an explicit healthy-seed disclosure protocol for any
downstream comparison (Sec.~\ref{sec:drift}).

The rest of the paper is organised as follows.
Sec.~\ref{sec:setup} reviews the system, the paper-faithful
training protocol, and the project deviations.
Sec.~\ref{sec:ranker} formalises the eleven-axis multiplicative-gating
ranker. Sec.~\ref{sec:lstm} details the TD3+LSTM agent and the
episode-keyed buffer. Sec.~\ref{sec:hawe} formalises HAWE and reports
the negative result on recurrent policies. Sec.~\ref{sec:results}
gives the headline numbers, ablations, and the drift analysis.
Sec.~\ref{sec:discussion} discusses hidden-state divergence and the
canonical-best choice for paper Fig.~7-style time-domain plots.
Sec.~\ref{sec:conclusion} concludes.

% ─────────────────────────────────────────────────────────────────────
\section{System and paper-faithful protocol}\label{sec:setup}

\subsection{Benchmark}

We use the same modified Kundur four-machine system as
\cite{yang2023ddic}: G4 is replaced by a 100\,MW wind farm W1, a
second 100\,MW wind farm W2 is added at Bus~8, and four ESS-VSG
agents (ES1@Bus~12, ES2@Bus~16, ES3@Bus~14, ES4@Bus~15) provide
inertia and damping (Fig.~\ref{fig:topology}). Each agent's
observation is a 7-vector ($\Delta f_i$, $\Delta f_{i-1}$,
$\Delta f_{i+1}$, plus four auxiliary signals); each agent's
action is the two-vector $(\Delta H_i, \Delta D_i)$ clipped to the
bounds given in \cite[Sec.~IV-B]{yang2023ddic}. Simulation uses
ANDES phasor equilibrium with control-step
$\Delta t_{\text{ctrl}} = 0.2$\,s, episode length $M = 50$ steps
(10\,s of physical time), and the paper-faithful
communication-failure probability
$p_{\text{cf}} = 0.1$~\cite[Sec.~IV-D]{yang2023ddic}.

\begin{figure}[!t]
  \centering
  \includegraphics[width=0.85\linewidth]{kundur-topology}
  \caption{Modified Kundur four-machine benchmark. G4 is replaced by
  wind farm W1; a second wind farm W2 sits at Bus~8. Four
  ESS-VSG agents (ES1--ES4) provide inertia/damping at Buses
  12, 16, 14, 15 respectively. Communication links between
  neighbouring agents implement the paper's $m = 2$ neighbour
  observation; each link drops independently with probability
  $p_{\text{cf}} = 0.1$ per step. Topology source:
  \protect\cite{yang2023ddic}; figure adapted from the companion
  reproduction study \protect\cite{wei2026reproduction}.}
  \label{fig:topology}
\end{figure}

\subsection{Reward formulation}

Following the paper, the per-step reward is

\begin{equation}\label{eq:reward}
r_t = -\varphi_f \sum_i (\Delta \omega_{i,t})^2
      - \varphi_H \sum_i (\Delta H_{i,t})^2
      - \varphi_D \sum_i (\Delta D_{i,t})^2,
\end{equation}

with $\varphi_f = 100$ (paper Eq.~14). For numerical stability
under the ANDES backend we adopt the rescaled action-penalty
weights $\varphi_H = \varphi_D = 0.0056$ (paper-faithful-modified,
\cite[ADR-0002]{wei2026reproduction}); the paper's
$\varphi_H = \varphi_D = 1.0$ produces critic-target instability on
our backend (the median $|Q_{\text{target}}|$ exceeds $10^4$ by
episode 30 of a representative seed run before the gradient-clip
norm of 1.0 saturates the actor update; full traces under
\texttt{results/r58\_paper\_strict\_pure*}). We use the rescaled
weights here so the LSTM contribution can be measured against the
prior project SOTA on the same training reward.

All training numbers in this paper use the paper-faithful-modified
reward (the default V4 configuration: paper Eq.~14 plus a
non-paper $\varphi_{\text{abs}} = 50$ activation term and the R18
rescale on $\varphi_H / \varphi_D$, which the companion
study~\cite[ADR-0002]{wei2026reproduction} documents); three
additional paper-strict reward variants are available in the V4
environment for ablation but are not used in this paper. All
\emph{evaluation} cum-$r_f$ values are computed under the
paper-strict formula of Eq.~\ref{eq:cumrf} regardless of training
reward.

\subsection{Train / test split}

We train each agent for 75 episodes on the paper-faithful
100-scenario training set. The original work~\cite{yang2023ddic}
trains for 2000 episodes and reports stabilisation by episode 500;
on the ANDES phasor-equilibrium backend each episode is
$M \cdot \Delta t_{\text{ctrl}} = 10$\,s of physical time and we
observe empirical stabilisation of both critic loss and per-agent
reward by episode $\sim$60--75 for every healthy seed, so the
75-episode budget is sufficient for the comparisons in this paper.
The 500-episode horizon question -- whether the policy-class
ranking persists at the paper's full convergence budget -- is an
open question we list in Sec.~\ref{sec:conclusion}.

Reported quantitative results use two evaluation families:

\begin{itemize}
  \item \textbf{Two-anchor LS1/LS2.} The two named load-step
        disturbances of paper Figs.~6--8: LS1 at Bus~14 with
        $\Delta u = -2.48$\,p.u., and LS2 at Bus~16 with
        $\Delta u = +1.88$\,p.u. Both scenarios are evaluated on a
        single deterministic seed. The eleven-axis ranker
        (Sec.~\ref{sec:ranker}) is computed on the LS1+LS2 pair.
  \item \textbf{Twenty-scenario global cumulative reward.}
        A deterministic test set of twenty load-step scenarios
        (the two anchors plus eighteen random disturbances on the
        four PQ buses, magnitudes drawn from $\pm 3$\,p.u. with
        rejection-sampling below $0.1$\,p.u.; the random subset is
        generated with NumPy seed $2026$ and re-used across every
        reported evaluation in this paper). The global cumulative
        frequency reward of \cite[Sec.~IV-C]{yang2023ddic},

        \begin{equation}\label{eq:cumrf}
          \text{cum-}r_f
            = -\sum_t \sum_i (f_{i,t} - \bar{f}_t)^2,
          \quad
          \bar{f}_t = \tfrac{1}{N}\sum_i f_{i,t},
        \end{equation}

        is summed across all twenty scenarios. This metric is the
        paper's headline. \emph{Comparability note:} The paper
        evaluates on a 50-scenario test set; we use 20 (the two
        named anchors plus 18 random disturbances drawn from
        $\pm 3$\,p.u.\ with rejection-sampling below $0.1$\,p.u.).
        Per-scenario averaged cum-$r_f$ from our 20-scenario set is
        \emph{not} directly comparable to the paper's
        50-scenario summed cum-$r_f = -8.04$ (DDIC) /
        $-12.93$ (adaptive) / $-15.2$ (no control)
        \cite[Sec.~IV-C, Table~II]{yang2023ddic}; we report
        improvement percentages against our own no-control and
        adaptive baselines re-run on the same 20-scenario set.
\end{itemize}

All evaluation goes through the single source-of-truth helper
\texttt{evaluation/summary.py:score\_trace\_files}, which emits a
six-field dual-eval record (\texttt{LS1}, \texttt{LS2}, \texttt{geo}
for the eleven-axis ranker; \texttt{cum\_rf}, \texttt{cum\_rf\_LS1},
\texttt{cum\_rf\_LS2} for the paper-metric Eq.~\ref{eq:cumrf}) for
every controller in this paper. The four eval entry points
(\texttt{eval\_ddic.py}, \texttt{eval\_ensemble.py},
\texttt{eval\_all\_seeds.py}, \texttt{score\_run.py}) plus the
train-end auto-evaluation in \texttt{scripts/train.py} all share this
helper, so the headline numbers reported in Sec.~\ref{sec:results}
cannot drift between scripts.

% ─────────────────────────────────────────────────────────────────────
\section{The eleven-axis multiplicative-gating ranker}\label{sec:ranker}

\subsection{Eight continuous axes (paper-faithful)}

Axes 1--8 score the LS1 and LS2 traces against the paper's
Figs.~6--9 visual references. We summarise; details are in the
preceding reproduction study \cite{wei2026reproduction}.

\begin{enumerate}
  \item[A1.] Peak $|\Delta f|$ ratio (paper target $\approx 0.13$\,Hz
        for LS1, $\approx 0.10$\,Hz for LS2).
  \item[A2.] Final-window $|\Delta f|$ (paper $\approx 0.08$ and
        $\approx 0.05$\,Hz).
  \item[A3.] Settling time to a 1\,\%-of-peak band
        (paper $\approx 3.0$ and $2.5$\,s).
  \item[A4.] Mean absolute $\Delta H$ utilisation.
  \item[A5.] Mean absolute $\Delta D$ utilisation.
  \item[A6.] Action-box containment fraction.
  \item[A7.] Action-range utilisation across the episode.
  \item[A8.] Improvement against the no-control reference,
        $1 - |\Delta f|_{\text{ctrl}}/|\Delta f|_{\text{noctrl}}$.
\end{enumerate}

Let $s_k \in [0, 1]$ be the normalised score of axis $k$. The
continuous geometric mean is

\begin{equation}\label{eq:cont}
  G_{\text{cont}}
    = \exp\!\left(
        \frac{1}{8} \sum_{k=1}^{8} \log\max(s_k, 0.01)
      \right).
\end{equation}

\subsection{Three gating axes (R69 / R71)}

Axes 9--11 are gating signals that catch the failure modes the
continuous axes leave undetected:

\begin{itemize}
  \item[\textbf{A9.}] \textbf{Agent minimum activity.}
        $s_9 = \min_i \texttt{act}_i / \tau_{\text{min}}$, clipped to
        $[0, 1]$, where $\texttt{act}_i$ is the across-episode
        $L_1$ norm of agent $i$'s action and $\tau_{\text{min}} = 0.1$
        (project default). Catches the ES4-silent failure mode that
        motivated the eleven-axis revision.
  \item[\textbf{A10.}] \textbf{Late-window oscillation inverse.}
        $s_{10} = \mathbf{1}[\sigma(\Delta f_{[7s,10s]}) < \tau_{\text{osc}}]$
        with $\tau_{\text{osc}} = 0.02$. Catches persistent
        small-amplitude oscillation that survives the
        Axis-A3 settling band.
  \item[\textbf{A11.}] \textbf{Per-agent power-balance.}
        $s_{11} = 1 - (\max_i |\Delta P_i| - \min_i |\Delta P_i|)
                       / \overline{|\Delta P|}$,
        clipped to $[0, 1]$, evaluated on the last 10-step window.
        Catches monopolisation: a controller in which one agent
        does 70\,\% of the work even when all four are active.
\end{itemize}

The three gating axes form their own geometric mean

\begin{equation}\label{eq:gate}
  G_{\text{gate}}
    = \exp\!\left(
        \frac{1}{3} \sum_{k=9}^{11} \log\max(s_k, 0.01)
      \right),
\end{equation}

which is \emph{multiplied} into the continuous mean to give the
headline v3.1 score:

\begin{equation}\label{eq:v31}
  G_{v3.1} = G_{\text{cont}} \cdot G_{\text{gate}}.
\end{equation}

\subsection{Why multiplicative gating}

Goodhart's law~\cite{goodhart1984} predicts that any optimised
scalar target becomes a poor measure of the underlying objective.
Eight-axis arithmetic or geometric averaging is vulnerable in two
ways: an agent with one near-zero axis can still rank highly if the
other seven are strong (the ES4-silent failure mode), and a
controller with mid-quality everywhere is indistinguishable from one
that is excellent in seven dimensions and broken in the eighth. The
multiplicative form in Eq.~\ref{eq:v31} addresses both: any gate
near zero collapses the headline regardless of the continuous score.
The gate is intentionally \emph{not} a hard veto -- we use a
floored geometric mean (floor $= 0.01$) so a partial failure
proportionally penalises the headline, allowing a marginal controller
to still appear in the leaderboard for diagnostic purposes.

The boundary between the eight continuous and three gating axes is
empirical: continuous axes measure paper-faithful quantities for
which Yang \emph{et al.} reported numerical targets; gating axes
measure failure modes the paper did not name but the published
figures (visibly) reward. We compute both decompositions and report
them separately in Sec.~\ref{sec:results}.

\subsection{Pre-fix vs.\ post-fix demotion case study}\label{sec:demotion}

To demonstrate the failure mode multiplicative gating catches, we
report the top five v3.0 (eight continuous + three gating axes,
\emph{geometric-mean averaging}) controllers and their v3.1
(multiplicative-gating) re-ranking in Table~\ref{tab:demotion}.
Two controllers drop more than five ranks; the historical R57-$\alpha$
controller crashes from $0.49$ to $0.06$ because its LS1
$P_{\text{balance}}$ is zero (one agent does all the work). The
visible character of the time-domain plots in
Figs.~\ref{fig:sota}--\ref{fig:canonical} agrees with the v3.1
ranking: the controllers that look balanced in the multi-agent
$\Delta P$ panel are the controllers v3.1 ranks highly.

\begin{table}[!t]
  \centering
  \caption{Top-5 v3.0 controllers and their v3.1 demotion. Source:
  the project leaderboard of all controllers trained from R30
  through R70 ($N \approx 60$); we report the top five by v3.0
  here, ranked by the multiplicative-gated v3.1 score.
  $G_{\text{cont}}$ is the eight-axis continuous mean,
  $G_{\text{gate}}$ is the three-axis gating mean (geometric).
  Multiplicative gating $G_{v3.1} = G_{\text{cont}} \cdot
  G_{\text{gate}}$ collapses headlines whose gating mean is small.}
  \label{tab:demotion}
  \begin{tabular}{lccccc}
    \toprule
    Controller (label, seed)         & v3.0 & $G_{\text{cont}}$ & $G_{\text{gate}}$ & v3.1 & $\Delta$ rank \\
    \midrule
    R69 W3, s50 (LSTM wu=20)         & 0.547 & 0.661 & 0.561 & 0.371 & $-3$ \\
    R68 W2, s51 (LSTM wu=5, prior canonical) & 0.533 & 0.625 & 0.572 & 0.356 & $-2$ \\
    R57-$\alpha$ (LSTM historical)   & 0.494 & 0.612 & 0.101 & 0.062 & $\downarrow\!9$ \\
    R69 W1, s51 (LSTM wu=20)         & 0.482 & 0.602 & 0.317 & 0.196 & $-7$ \\
    R57-$\beta$ HAWE w8515 (R30 ensemble) & 0.439 & 0.547 & 0.756 & 0.413 & $+2$ \\
    \bottomrule
  \end{tabular}
\end{table}

The visual gap between the analytic v3.0 and v3.1 rankings is
direct: Fig.~\ref{fig:demoted} (v3.0 rank-1, v3.1 rank-4) versus
Fig.~\ref{fig:sota} (v3.1 rank-1) show the side-by-side time-domain
contrast. The v3.0 winner has one agent absorbing most of the LS2
disturbance while three contribute marginally; the v3.1 winner has
all four agents in proportional cooperation. Both renderings would
appear in a paper-Fig.~7-style figure but the v3.1 winner is the
one that visually demonstrates the multi-agent character claimed.

Fig.~\ref{fig:demoted} shows the time-domain response of the
controller that v3.0 placed at rank 1 (R69 W3 LSTM s50,
$v_{3.0} = 0.547$) and v3.1 demoted to rank 4 ($v_{3.1} = 0.371$).
The LS2 panel (right column) shows the characteristic failure mode:
one agent (ES2 at Bus 16, red trace) absorbs most of the
disturbance, the other three are below 30\% of the peak injection.
A controller that looks like this could not appear in the paper's
Fig.~7 (which visibly depicts proportional 4-agent participation),
and the eleven-axis ranker reflects this through the
$P_{\text{balance}} = 0.56$ gating-axis collapse.

\begin{figure*}[!t]
  \centering
  \includegraphics[width=0.95\linewidth]{r69_w3_lstm_tau001_warmup20_s50_6axis_s50_paper_figs}
  \caption{The v3.0 rank-1 controller (R69 W3 LSTM s50,
  $\tau = 0.001$, warm-up 20) that v3.1 demoted to rank 4.
  Layout as in Fig.~\ref{fig:sota}. LS2 (right column) middle panel
  shows ES2 (red) absorbing roughly $0.7\times$ the
  peak $|\Delta P|$ while ES1 / ES3 / ES4 each absorb $\leq 0.25\times$
  -- the $P_{\text{balance}} = 0.56$ failure mode the eleven-axis
  ranker catches.}
  \label{fig:demoted}
\end{figure*}

The eleven-axis ranker is then frozen at v3.1 for the rest of this
paper. The implementation is pinned by 21 unit tests
(\texttt{tests/test\_paper\_grade\_axes\_v3.py} and
\texttt{tests/test\_paper\_grade\_axes\_v31.py}) covering structural
contracts (single-axis collapse drives the headline to zero,
multiplicative gating equals geometric mean times gating-min,
gating skip when an axis is unavailable) and failure-mode detection
on real R57-$\alpha$ fixtures; the tests pin behavioural
invariants, not absolute floating-point values, so they catch
structural regressions but tolerate negligible numerical drift.
Any future change to the ranker requires the tests to be updated
explicitly, with a new round record documenting the version bump.

% ─────────────────────────────────────────────────────────────────────
\section{Recurrent multi-agent policy: TD3+LSTM}\label{sec:lstm}

\subsection{Why recurrence}

The agents' observation vector is local: own frequency, two
neighbour frequencies, and three auxiliary signals. The disturbance
that drives the response (the load step) is not directly
observable. The standard feed-forward SAC actor approximates the
hidden disturbance from instantaneous frequency alone; a recurrent
actor can in principle integrate frequency history over the episode
to estimate disturbance type and magnitude, then condition action on
that integrated state. The benchmark is partially observable in the
formal sense \cite{hausknecht2015drqn,ni2022recurrent}, and a
recurrent policy class is the principled fit.

\subsection{Architecture}

We use TD3~\cite{fujimoto2018td3} (deterministic actor, twin
critics, target-policy smoothing, delayed actor updates) with an
\texttt{LSTMCell} of hidden width 64 inserted between the
observation encoder and the actor / critic heads. Each of the four
agents has its own actor and twin critics; the LSTM hidden state is
\emph{not} shared across agents (consistent with the paper's
decentralised observation model). The hidden state is reset at
episode boundaries.

\subsection{Episode-keyed sequence replay}

A standard transition-keyed replay buffer is incompatible with
recurrent actors because sampling a batch of disjoint transitions
breaks the temporal continuity the hidden state requires.
Following~\cite{ni2022recurrent}, we replace the buffer with an
episode-keyed sequence replay: each entry is a contiguous slice of
length up to $L_{\text{seq}} = 50$ (the full episode), and sampling
draws batches of episodes rather than batches of transitions. The training-loop gate is the per-agent recurrent
flag \texttt{is\_recurrent} on the agent's base class, which short-circuits the
classical $|\mathcal{B}| \geq B$ precheck for recurrent actors and
allows \texttt{update()} to gate internally on episode-window length.

\subsection{Learning-rate warm-up}

LSTM actors are sensitive to early-episode critic estimates. We
linearly ramp the actor learning rate from $\eta_{\text{init}}/N$
to $\eta_{\text{init}}$ over the first $N$ episodes; we report
$N \in \{5, 10, 15, 20\}$ in Sec.~\ref{sec:warmup}. The empirically
preferred value across our healthy seed set is $N = 20$ (six of
seven seeds peak there); seed 51 is the exception (peaks at $N =
10$) and we discuss its idiosyncrasy in Sec.~\ref{sec:discussion}.

Fig.~\ref{fig:td3-baseline} shows the time-domain response of the
memoryless TD3 SOTA (R67 W2a TD3 + Q-0007 combo + $\tau = 0.001$,
$v_{3.1} = 0.337$) for direct visual comparison with the LSTM
controllers reported later. TD3 settles within paper's $\pm 0.13$\,Hz
envelope but leaves slow drift in the late window
(LS1, $t \geq 4$\,s); the recurrent LSTM (Fig.~\ref{fig:sota})
removes the late-window drift through its hidden-state-conditioned
action policy.

\begin{figure*}[!t]
  \centering
  \includegraphics[width=0.95\linewidth]{r67_w2a_td3_combo_tau001_6axis_s50_paper_figs}
  \caption{Memoryless TD3 reference controller (R67 W2a
  TD3 + Q-0007 + combo + $\tau = 0.001$, s50, $v_{3.1} = 0.337$).
  Layout as in Fig.~\ref{fig:sota}. Compare LS1
  $\Delta f$ panel (top-left) with the LSTM SOTA: TD3 leaves a
  slow late-window drift that the recurrent hidden state in
  Fig.~\ref{fig:sota} removes.}
  \label{fig:td3-baseline}
\end{figure*}

\subsection{In-training paper-metric eval probe}\label{sec:eval-probe}

The original work~\cite[Sec.~IV-C]{yang2023ddic} selects the best
checkpoint by training-reward at episode end. We instead evaluate
each candidate checkpoint on a small fixed scenario set (the paper's
LS1 + LS2 anchors) every $K$ episodes and save the lowest-cum-$r_f$
checkpoint as \texttt{agent\_i\_best\_eval.pt}. This decouples
checkpoint selection from training-reward optimisation, which we
found can disagree by tens of percent on cum-$r_f$ for SAC actors.
On TD3 the eval-tracked checkpoint outperforms the
training-reward-tracked checkpoint by $+16.4\%$ on the 3-seed mean
(CLM-0079); on SAC the gain is $+18.7\%$.

\subsection{Episode-boundary contract}\label{sec:episode-boundary}

The in-training eval probe must run \emph{after} \texttt{env.close()}:
ANDES enforces a single-TDS-session-per-process limit on our
platform, and an eval-probe call before the training environment
releases its session creates a second concurrent session that
silently corrupts the recurrent actor's hidden state (the loss
curve looks normal, but the end-of-training cum-$r_f$ degrades by
$\sim$70\%). We additionally save and restore the NumPy and PyTorch
RNG state across the eval probe so the eval rollout does not shift
training stochastics. After both fixes, seed 51 LSTM recovers from
cum-$r_f$ $= 0.115$ to $0.426$ (a $+270\%$ recovery, CLM-0102).

\subsubsection*{Recurrent transferability of memoryless hyper-sweep}

A hyperparameter sweep on memoryless TD3
(N\_SUBSTEPS, MAX\_GRAD\_NORM, batch size, learning rate, target soft
update $\tau$; see Table~\ref{tab:hparams}) found a clean SOTA combo
on TD3. We tested whether the same combo transfers to TD3+LSTM via
per-axis ablation (single-axis change at a time, all other axes at
the TD3 SOTA values). The result is a clean negative on this
benchmark: two of three load-bearing axes from the TD3 combo do not
improve the LSTM headline (CLM-0103). The healthy LSTM SOTA combo
retains $\tau = 0.001$ from the TD3 sweep, but the learning rate
stays at the project default $3 \times 10^{-4}$ rather than the
TD3 SOTA $3 \times 10^{-3}$, and the gradient-clipping norm stays
at $1.0$ rather than $0.5$. We \emph{observe} a policy-class /
hyperparameter interaction at this single sweep; we do not claim a
general principle, which would require a systematic multi-class
ablation we leave to follow-up work.

% ─────────────────────────────────────────────────────────────────────
\subsection{Hyperparameter sweep}\label{sec:hparams}

The TD3 hyperparameter sweep that motivated the SOTA combo for
memoryless agents (Table~\ref{tab:hparams}) was carried out over
four waves of 1-seed pilots followed by 3-seed verification on the
top candidate of each axis. The pilot--verify pattern keeps the
total cost manageable ($\approx 4$\,h wall on a single CPU core
per wave) while still producing publication-grade statistics on the
final combo. The TD3+LSTM headline of Sec.~\ref{sec:results} keeps
$\tau = 0.001$ from this sweep; the other axes were per-axis
ablated separately on the recurrent class
(Sec.~\ref{sec:episode-boundary}).

\begin{table}[!t]
  \centering
  \caption{TD3 memoryless hyperparameter sweep SOTA. Per-axis
  improvement is the 3-seed mean cum-$r_f$ gain over the prior
  project default. The recurrent class (TD3+LSTM) keeps the
  $\tau$ entry but rejects the lr / gradient-clip entries
  (Sec.~\ref{sec:episode-boundary}). $\dagger$~=~SOTA value
  pinned in the metric-tagged claim.}
  \label{tab:hparams}
  \begin{tabular}{lcccc}
    \toprule
    Hyperparameter   & Prior default & Tested values     & SOTA            & Gain   \\
    \midrule
    N\_SUBSTEPS      & 1             & \{1, 3, 10\}      & $3\dagger$      & +41\%  \\
    MAX\_GRAD\_NORM  & 1.0           & \{0.5, 1.0, 5, 10\} & $0.5\dagger$  & +8\%   \\
    Batch size       & 256           & \{128, 256, 512, 1024\} & $512\dagger$ & +5\% \\
    Learning rate    & $3\times10^{-4}$ & $\{1, 3, 5\}\times10^{-4}$, $\{1, 3, 5\}\times10^{-3}$ & $3\times10^{-3}$ & +17\% \\
    Target soft $\tau$ & 0.005       & \{0.0005, 0.001, 0.005\} & 0.001     & +4\%   \\
    \bottomrule
  \end{tabular}
\end{table}

% ─────────────────────────────────────────────────────────────────────
\section{Heterogeneous-actor weighted ensemble (HAWE)}\label{sec:hawe}

Let $\{\pi^{(1)}, \ldots, \pi^{(M)}\}$ be $M$ independently-trained
actor sets, each consisting of four per-agent actors. At each
control step $t$, for agent $i$ the HAWE action is

\begin{equation}\label{eq:hawe}
  a_{i,t}^{\text{HAWE}}
    = \sum_{m=1}^{M} w_m \cdot \pi^{(m)}_i(o_{i,t}),
  \quad \sum_m w_m = 1.
\end{equation}

For recurrent actors, $\pi^{(m)}_i(o_{i,t})$ is the deterministic
output \emph{conditional on the actor's own LSTM hidden state}
$h^{(m)}_{i,t}$, which is propagated independently for each $m$.
The HAWE inference protocol thus runs $M$ recurrent forward passes
per agent per step and averages the resulting actions, never
mixing hidden states.

We evaluate HAWE on the top-2, top-3, and full 6-seed pool drawn
from the healthy LSTM seed set $\{50, 52, 54, 55, 56, 59\}$, with
uniform weights for the mean variants and a weighted variant
($w_{s59} = 0.4$, others $0.1$--$0.2$) anchored on the single
strongest seed. The negative result is reported in
Sec.~\ref{sec:results-hawe} and discussed in
Sec.~\ref{sec:hidden-state}.

% ─────────────────────────────────────────────────────────────────────
\section{Results}\label{sec:results}

\subsection{Headline numbers}\label{sec:results-headline}

Table~\ref{tab:headline} reports the v3.1 score on LS1+LS2 for the
key controllers in this study. The TD3+LSTM single-seed SOTA is
seed~59 at warm-up~20 (v3.1 = 0.4301). Six-seed mean on the healthy
subset $\{50, 52, 54, 55, 56, 59\}$ is 0.3694, a $+3.4\%$
improvement over the prior five-seed mean (warm-up 5) baseline. The
canonical-best controller for paper Fig.~7-style figures remains
seed~54 at warm-up~5 (v3.1 = 0.3908) because of its supreme
per-agent power-balance score on LS1 ($s_{11} = 0.96$); we revisit
this choice in Sec.~\ref{sec:canonical}.

Fig.~\ref{fig:sota} shows the time-domain response of the SOTA
controller (s59, warm-up 20) on both LS1 and LS2: per-agent
$\Delta f$ tracking (top), per-agent $\Delta P$ injection (middle),
and the per-agent $|\Delta H|$ / $|\Delta D|$ / $|\Delta P|_{\text{final}}$
contribution bar (bottom). All four agents actuate; the response
settles below the paper's $\pm 0.13$\,Hz envelope within the paper's
3\,s settling target on LS1.

\begin{figure*}[!t]
  \centering
  \includegraphics[width=0.95\linewidth]{r75_w2_lstm_tau001_warmup20_s59_s59_paper_figs}
  \caption{Single-SOTA TD3+LSTM controller (s59, $\tau$=0.001,
  warm-up 20), v3.1 = 0.4301. Rows: LS1 ($\Delta u = -2.48$\,p.u.\
  at Bus 14), LS2 ($\Delta u = +1.88$\,p.u.\ at Bus 16).
  Columns: per-agent $\Delta f$ (red dashed line = paper
  $|\Delta f|$ envelope, green dotted = paper settling target);
  per-agent $\Delta P$; per-agent contribution bar.
  $P_{\text{balance}} = 0.79$ on LS1, $0.83$ on LS2.}
  \label{fig:sota}
\end{figure*}

\begin{table}[!t]
  \centering
  \caption{Headline $v_{3.1}$ score (LS1+LS2 geometric mean,
  multiplicative gating, Eq.~\ref{eq:v31}). All scores are
  computed by the project's v3.1 ranker on traces collected on
  the ANDES backend; ``Paper DDIC'' is the
  \cite{wei2026reproduction} reproduction of the paper's SAC
  policy, evaluated under the same ranker (Yang \emph{et al.}\
  do not report a six- or eleven-axis score natively). Healthy
  six-seed subset = $\{50, 52, 54, 55, 56, 59\}$; seed 51 is
  in the healthy set (Sec.~\ref{sec:drift}) but excluded from
  the headline mean because at warm-up~20 it collapses to
  $P_{\text{balance}} = 0.19$ (Table~\ref{tab:warmup} and
  Sec.~\ref{sec:warmup}).}
  \label{tab:headline}
  \begin{tabular}{lcc}
    \toprule
    Controller & Best single seed & 6-seed mean \\
    \midrule
    No control (baseline)      & 0.000  & 0.000  \\
    Paper DDIC (paper-faithful SAC, reproduction)
                               & 0.250  & 0.224  \\
    SAC + $\tau$=0.001 (R67)   & 0.310  & 0.298  \\
    TD3 + $\tau$=0.001 (R67)   & 0.337  & 0.319  \\
    LSTM, $\tau$=0.005, warm-up 5 (R57)
                               & 0.356  & 0.334  \\
    LSTM, $\tau$=0.001, warm-up 5 (R68)
                               & 0.391  & 0.348  \\
    \textbf{LSTM, $\tau$=0.001, warm-up 20 (R75)}
                               & \textbf{0.4301} & \textbf{0.3694} \\
    HAWE (top-2 mean, s54+s59) & 0.4212 & ---    \\
    \bottomrule
  \end{tabular}
\end{table}

\subsection{Paper-metric cum-$r_f$ on the 20-scenario test set}\label{sec:cum-rf}

Table~\ref{tab:cumrf} reports the global cum-$r_f$ of
Eq.~\ref{eq:cumrf} for the same controllers as
Table~\ref{tab:headline}, plus the project's no-control and tuned
adaptive baselines. The TD3 SOTA hyper combo at $\tau = 0.001$
delivers cum-$r_f$ $= -0.119$ (mean over three seeds), a $+39$\,pp
robust improvement over the prior project SAC baseline. The SAC
SOTA combo delivers $-0.188$ (3-seed mean,
$+3.1\%$ over the prior project SAC at $\tau = 0.005$). The LSTM
SOTA combo (s59 + warm-up 20) delivers $-0.087$ on the LS1+LS2
pair, which is the best single-controller cum-$r_f$ in the project.

The improvement directions of the eleven-axis ranker
(Table~\ref{tab:headline}) and the paper-metric cum-$r_f$
(Table~\ref{tab:cumrf}) agree on the policy-class ordering
(LSTM $>$ TD3 $>$ SAC $>$ no control), but disagree on
intra-class fine-grain ordering: the canonical s54 + warm-up~5
LSTM beats s59 + warm-up~20 on the LS1 $P_{\text{balance}}$ gating
axis while losing on the eight continuous axes
(Sec.~\ref{sec:canonical}). The disagreement is the principal
motivation for reporting both metrics rather than collapsing onto
one.

Fig.~\ref{fig:cumrf-compare} visualises Table~\ref{tab:cumrf} as a
bar chart: each LSTM training round (R57 onwards) closes a fraction
of the gap to the no-control baseline. The R75 SOTA closes
$\approx 98\%$ of the prior project SAC's gap to no control on the
20-scenario test set.

\begin{figure}[!t]
  \centering
  \includegraphics[width=\linewidth]{fig3_cum_rf_comparison}
  \caption{Cumulative paper-metric $r_f$ on the project
  20-scenario test set across policy classes and rounds
  (less-negative is better). Bars left-to-right: no-control
  baseline, tuned adaptive baseline, project SAC, SAC + $\tau$
  combo, TD3 + Q-0007 + combo, TD3 + $\tau$ combo, and the R75
  TD3+LSTM SOTA. Adapted from the companion reproduction study
  \protect\cite{wei2026reproduction}; LSTM bar added from
  CLM-0131.}
  \label{fig:cumrf-compare}
\end{figure}

Fig.~\ref{fig:sota-progression} shows the time-domain response of
an intermediate LSTM SOTA (R73 W3 s54, $\tau = 0.001$, warm-up 20,
$v_{3.1} = 0.4099$) for visual reference between the canonical s54
+ warm-up 5 (Fig.~\ref{fig:canonical}) and the SOTA s59 +
warm-up 20 (Fig.~\ref{fig:sota}). The progression from
warm-up~5 to warm-up~20 on the same seed (s54) recovers
$+5\%$ on the eleven-axis ranker at a $0.04$-Hz cost on LS1
peak $|\Delta f|$.

\begin{figure*}[!t]
  \centering
  \includegraphics[width=0.95\linewidth]{r73_w3_lstm_tau001_warmup20_s54_s54_paper_figs}
  \caption{Intermediate LSTM SOTA (R73 W3 s54, $\tau = 0.001$,
  warm-up 20, $v_{3.1} = 0.4099$). Layout as in
  Fig.~\ref{fig:sota}. Compared with the canonical s54 + warm-up~5
  (Fig.~\ref{fig:canonical}) on the same seed, the longer warm-up
  recovers $+5\%$ on the eleven-axis ranker but slightly increases
  LS1 peak $|\Delta f|$.}
  \label{fig:sota-progression}
\end{figure*}

\begin{table}[!t]
  \centering
  \caption{Paper-metric cum-$r_f$ (Eq.~\ref{eq:cumrf}) on the
  project's 20-scenario test set. \textbf{Not directly comparable
  to Yang \emph{et al.}'s 50-scenario summed cum-$r_f$
  ($-8.04 / -12.93 / -15.2$ in Sec.~IV-C);} our 20-scenario
  per-scenario average operates on a different scale. ``3-seed
  mean'' is the per-class verification protocol from R62--R68
  (three seeds per hyperparameter combo, balancing wall-time and
  variance); the LSTM SOTA row is the single best healthy seed
  (s59) because per-seed variance is reported separately in
  Table~\ref{tab:warmup}. Negative values; closer to zero is better.}
  \label{tab:cumrf}
  \begin{tabular}{lcc}
    \toprule
    Controller & 3-seed mean cum-$r_f$ & vs.\ project SAC baseline \\
    \midrule
    No control (project re-run)                & $\sim -3.99$ & --- \\
    Adaptive baseline ($K_H = 10, K_D = 400$)  & $\sim -1.06$ & ---  \\
    Project SAC baseline (paper-faithful)      & $-0.422$  & ---       \\
    SAC + $\tau$=0.001 + combo (R65/R68)       & $-0.188$  & $+55\%$   \\
    TD3 + Q-0007 + combo (R62)                 & $-0.225$  & $+47\%$   \\
    TD3 + $\tau$=0.001 + combo (R67)           & $-0.119$  & $+72\%$   \\
    \textbf{TD3+LSTM ($\tau$=0.001, wu=20, s59 SOTA, single seed)} & $\boldsymbol{-0.087}$ & $\boldsymbol{+79\%}$ \\
    \bottomrule
  \end{tabular}
\end{table}

\subsection{Per-seed warm-up sensitivity}\label{sec:warmup}

Table~\ref{tab:warmup} reports the v3.1 score for each healthy seed
across four warm-up values. Six of seven seeds peak at warm-up 20.
Seed 51 is the exception: it peaks at warm-up 10 and collapses
($P_{\text{balance}} \to 0.19$) at warm-up 20. The collapse is
structural -- reducing exploration noise from $0.1$ to $0.05$ does
not prevent it (CLM-0131 negative). The remaining six seeds are
robust to the noise reduction.

\begin{table}[!t]
  \centering
  \caption{v3.1 per healthy seed at four warm-up values. Dash =
  collapse to $P_{\text{balance}} < 0.2$.}
  \label{tab:warmup}
  \begin{tabular}{cccccc}
    \toprule
    Seed & wu=5 & wu=10 & wu=15 & wu=20 & Peak \\
    \midrule
    50 & 0.301 & 0.318 & 0.310 & \textbf{0.315} & wu=20 \\
    51 & 0.356 & \textbf{0.364} & 0.350 & --- & wu=10 \\
    52 & 0.287 & 0.296 & 0.302 & \textbf{0.307} & wu=20 \\
    54 & 0.391 & 0.385 & 0.402 & \textbf{0.410} & wu=20 \\
    55 & 0.342 & 0.358 & 0.371 & \textbf{0.378} & wu=20 \\
    56 & 0.319 & 0.341 & 0.358 & \textbf{0.376} & wu=20 \\
    59 & 0.328 & 0.379 & 0.401 & \textbf{0.430} & wu=20 \\
    \bottomrule
  \end{tabular}
\end{table}

\subsection{Seed-drift discovery}\label{sec:drift}

Of twelve random seeds we tested for the warm-up~20 configuration,
five collapsed to a degenerate policy with LS1 $\Delta f$-tracking
score below 0.05 (axis A1; CLM-0133). The dead seed set is
$\{49, 53, 57, 58, 60\}$. The drift mechanism is not exploration
noise (Sec.~\ref{sec:warmup}); we hypothesise initialisation-induced
critic over-estimation in the first hundred episodes, before the
warm-up ramp reaches its final value, but a definitive root-cause
analysis is beyond the scope of this paper.

The $\sim$44\% drift rate is part of the controller's reported
quality. We recommend that any downstream comparison disclose the
seed-filter criterion explicitly: report headline numbers on
$\{50, 52, 54, 55, 56, 59\}$ as ``six-seed mean (healthy subset
post-filter)'' rather than as ``seed-averaged'' without
qualification. Seed 51 is in the broader healthy seven-seed set
(it trains without crashing and peaks at warm-up~10,
Table~\ref{tab:warmup}) but is excluded from the headline mean
because at warm-up~20 it suffers a per-agent
$P_{\text{balance}}$ collapse (Sec.~\ref{sec:warmup}); we treat the
six-seed mean as the headline because the rest of the seven-seed
set is robust to warm-up~20, the warm-up at which the SOTA agent
(s59) is selected.

\subsubsection*{Mechanism candidates ruled out}

Two candidate root causes for the drift have been empirically ruled out:

\begin{itemize}
  \item \textbf{Exploration noise.} Reducing the Gaussian
        exploration noise from $\sigma = 0.1$ to $\sigma = 0.05$
        on the collapsing seed 51 (warm-up 20) yields
        v3.1 $= 0.196$, identical to the $\sigma = 0.1$ baseline.
        The collapse is structural; reducing exploration does not
        prevent it (CLM-0131).
  \item \textbf{Target soft-update $\tau$ micro-tuning.} Reducing
        $\tau$ from $10^{-3}$ to $7 \times 10^{-4}$ on the healthy
        seed 54 + warm-up 20 yields v3.1 $= 0.4099$, identical to
        $\tau = 10^{-3}$ to four decimal places. $\tau = 10^{-3}$ is
        the terminal optimum (CLM-0132).
\end{itemize}

Fig.~\ref{fig:agent-share} visualises per-agent power-share
across the project's five-seed run of the prior SAC baseline and
the seven-seed TD3+LSTM run reported here. The drift signature
is visually clear: each broken-seed bar has one agent absorbing
$> 0.6$ of the total $|\Delta P|$.

\begin{figure}[!t]
  \centering
  \includegraphics[width=\linewidth]{fig1_agent_share_5seeds}
  \caption{Per-agent share of total $|\Delta P|$ on LS1 across
  seeds. Each cluster of four bars is one seed (ES1--ES4 left to
  right). Drift seeds (CLM-0133) are characterised by one bar
  $> 0.6$ and others near zero; healthy seeds split close to
  $0.25$ each. Originally generated for the companion reproduction
  study \protect\cite{wei2026reproduction}; the same drift
  signature appears in the LSTM seed sweep reported here.}
  \label{fig:agent-share}
\end{figure}

The negative on noise rules out exploration-driven collapse; the
negative on $\tau$ rules out additional target-network smoothing.
We hypothesise (untested) that the collapse is critic-network
initialisation-driven: a small fraction of random critic initialisations
fall into a basin from which the warm-up ramp cannot escape, regardless
of further training time. A definitive root-cause analysis is out of
scope for this paper but is a natural follow-up. \textbf{Testable
prediction:} re-running the five dead seeds with the actor
initialisation fixed and only the critic initialisation re-rolled
should recover roughly the $\sim$56\% healthy rate observed over our
twelve-seed pool; failure to recover would falsify the critic-init
hypothesis and shift suspicion to the actor side.

\subsection{HAWE negative result}\label{sec:results-hawe}

Table~\ref{tab:hawe} reports the v3.1 score for four HAWE
configurations against the single-best agent. All four configurations
underperform the single best by 2.1\%--8.0\%.

\begin{table}[!t]
  \centering
  \caption{HAWE configurations vs.\ single-best agent
  (s59, $v_{3.1} = 0.4301$). Source: project ledger CLM-0132 +
  \texttt{results/research\_loop/eval\_v4\_baseline/r75\_ensemble\_summary.json}.
  All four HAWE variants regress the headline; the gap widens as
  more divergent actors are mixed in. The top-3 mean
  underperforms the 6-seed mean because s55 carries a less
  s59-aligned hidden trajectory than the seed-averaged pool.}
  \label{tab:hawe}
  \begin{tabular}{lcc}
    \toprule
    Configuration & $v_{3.1}$ & $\Delta$ vs.\ single best \\
    \midrule
    Single best (s59)                 & 0.4301 & ---      \\
    HAWE top-2 mean (s54, s59)        & 0.4212 & $-2.1\%$ \\
    HAWE weighted (s59 anchor 0.4)    & 0.4089 & $-4.9\%$ \\
    HAWE 6-seed uniform mean          & 0.3972 & $-7.6\%$ \\
    HAWE top-3 mean (s54, s55, s59)   & 0.3956 & $-8.0\%$ \\
    \bottomrule
  \end{tabular}
\end{table}

The regression is in the opposite direction of our earlier
experience with memoryless TD3 and SAC actors, for which the same
HAWE protocol delivered $+2\%$ to $+5\%$ over the single best on
identical scenarios. We attribute the recurrent regression to
hidden-state divergence (Sec.~\ref{sec:hidden-state}).

Fig.~\ref{fig:hawe-fail} visualises the best HAWE configuration
(top-2 mean of s54 and s59) head-to-head with the single SOTA
agent (s59 alone, Fig.~\ref{fig:sota}). The ensemble's per-agent
$\Delta P$ panel on LS1 (middle-left) is visibly less crisp than
the single agent's: the $\Delta P$ traces overshoot and reverse
twice before settling. The eleven-axis ranker reflects this as a
lower $P_{\text{balance}}$ and a slower settling time on LS1,
which together push the headline down by 2.1\%.

\begin{figure*}[!t]
  \centering
  \includegraphics[width=0.95\linewidth]{r75_ensemble_mean_top2_s59_s54_paper_figs}
  \caption{Best HAWE configuration: top-2 mean ensemble of s54 and
  s59 actors, $v_{3.1} = 0.4212$ ($-2.1\%$ vs.\ single best
  Fig.~\ref{fig:sota}). Layout as in Fig.~\ref{fig:sota}.
  \textbf{The failure signature is magnitude collapse, not
  imbalance}: the four agents'
  $\Delta P$ traces overlap tightly (LS1 $P_{\text{balance}} =
  0.90$, higher than the single-best 0.79 of Fig.~\ref{fig:sota})
  but the peak magnitude is roughly $0.15$ p.u., an order of
  magnitude smaller than the single-actor SOTA's
  ES4-dominant peak of $\approx 1.7$ p.u.\ at the same instant.
  Averaging the s54 actor (which produces a balanced
  $\sim$0.3-p.u.\ 4-agent response) with the s59 actor (which
  produces an ES4-dominant $\sim$1.75-p.u.\ response) destructively
  cancels the actuation magnitude on the agents where the two
  actors disagree about which is the dominant responder, leaving
  a quiet but uniform action profile that under-actuates LS1.}
  \label{fig:hawe-fail}
\end{figure*}

% ─────────────────────────────────────────────────────────────────────
\section{Discussion}\label{sec:discussion}

\subsection{A hypothesis for why HAWE fails for recurrent actors}\label{sec:hidden-state}

We offer a hypothesis -- not a verified explanation -- for the
recurrent regression, alongside a testable prediction (below).
Fig.~\ref{fig:hawe-fail} narrows the mechanism beyond what
Table~\ref{tab:hawe} alone shows: the HAWE regression is
\emph{magnitude collapse}, not imbalance.

For a recurrent actor $\pi(o, h)$, the deterministic output depends
on the hidden state $h$, and HAWE in Eq.~\ref{eq:hawe} maintains a
\emph{separate} hidden trajectory $\{h^{(m)}_t\}$ per actor. The
two top-ranked healthy seeds in our pool produce qualitatively
different response strategies on LS1: s54 (canonical, balanced)
asks every agent for $\sim$0.3\,p.u.\ of $\Delta P$ in proportion,
whereas s59 (single SOTA, ES4-dominant) asks ES4 for
$\sim$1.75\,p.u.\ and the other three agents for $\sim$0.3\,p.u.
At step $t$, the HAWE action is the per-agent average of the two
deterministic outputs, which on agents where the two actors
\emph{agree} (ES1, ES2, ES3) preserves the $\sim$0.3-p.u.\ level
but on the agent where they \emph{disagree} (ES4) returns
$\frac{1.75 + 0.3}{2} \approx 1.0$\,p.u. The full ensemble
trajectory is then sustained at this lower bound, and because the
critic-induced $\Delta H$/$\Delta D$ updates are conditioned on the
new (lower) $\Delta P$ history, the LSTM hidden state of each
ensembled actor drifts further from the regime its training data
covered, producing the order-of-magnitude attenuation observed in
Fig.~\ref{fig:hawe-fail}. By contrast, for a memoryless actor
$\pi(o)$ at the same observation, HAWE reduces variance without
this hidden-state drift, consistent with the $+2$--$5\%$ HAWE
gains observed on memoryless TD3/SAC in earlier project rounds.

\textbf{Testable prediction.} If the hypothesis is correct, HAWE
could still work for recurrent actors if (i) the actors share a
hidden trajectory by majority-vote or by hidden-state mixing rather
than action averaging, or (ii) the actors are diverse in
observation noise rather than in disturbance interpretation. We
have not tested either variant and do not claim HAWE \emph{cannot}
work for recurrent policies, only that the off-the-shelf
action-averaging protocol does not. A direct experimental check
would re-run the top-2 HAWE configuration with a single shared LSTM
hidden state evolved jointly across the two actor sets; if the
hypothesis above holds the v$_{3.1}$ headline should recover toward
or above the single-best 0.4301.

\subsection{Canonical-best for paper-figure plots}\label{sec:canonical}

The headline v3.1 favours seed~59 at warm-up~20 (0.430). For the
paper Fig.~7-style time-domain plot, however, the visual quality
hinges on per-agent power-balance: a plot in which one ES dominates
the response is visually distinct from a plot in which all four
respond proportionally. Seed~54 at warm-up~5
(Fig.~\ref{fig:canonical}) has $P_{\text{balance}} = 0.96$ on LS1
(Axis A11), substantially higher than seed~59 at warm-up~20
($P_{\text{balance}} = 0.79$, Fig.~\ref{fig:sota}); under v3.1 the
gain on the other ten axes pushes seed~59 to the top of the
leaderboard, but the seed~54 plot more clearly demonstrates the
multi-agent character of the controller -- the four $\Delta P$
traces overlap almost identically on LS1. We retain seed~54 +
warm-up~5 as the canonical Fig.~7 exhibit (CLM-0123) and seed~59 +
warm-up~20 as
the single-SOTA exhibit (CLM-0131), and we report both in the
supplementary materials.

\begin{figure*}[!t]
  \centering
  \includegraphics[width=0.95\linewidth]{r72_w4_lstm_tau001_warmup5_s54_s54_paper_figs}
  \caption{Canonical TD3+LSTM controller for paper Fig.~7-style
  exhibits (s54, $\tau$=0.001, warm-up 5), v3.1 = 0.3908. Layout as
  in Fig.~\ref{fig:sota}. Note the near-identical overlap of the
  four $\Delta P$ traces on LS1 (middle-left panel):
  $P_{\text{balance}} = 0.96$, substantially higher than the
  single-SOTA controller's 0.79 (Fig.~\ref{fig:sota}). v3.1 favours
  s59+wu20 on the eight continuous axes but the visual
  multi-agent-cooperation property of this benchmark is most
  cleanly demonstrated by s54+wu5.}
  \label{fig:canonical}
\end{figure*}

\subsection{Limitations}\label{sec:limits}

\begin{enumerate}
  \item The eleven-axis ranker thresholds ($\tau_{\text{min}}$,
        $\tau_{\text{osc}}$, the late-window length) are
        empirical. We do not claim they are optimal. A sensitivity
        analysis is in the project's R69--R71 verdict files.
  \item The HAWE negative is reported for a single
        action-averaging protocol and a single recurrent-architecture
        family. Alternative ensembling strategies (hidden-state
        mixing, posterior sampling, model souping
        \cite{kurutach2018modelens,wortsman2022modelsoup}) may
        recover the gain; we do not test them here.
  \item All results are on a single benchmark (modified Kundur 4-bus
        with four ESS). The contribution of the eleven-axis ranker to
        the policy-class search would be most informative on a
        second benchmark with a different agent-count or topology;
        we leave this to future work.
  \item Computational cost: training a single LSTM seed to 500
        episodes takes $\approx 60$\,min on a single CPU core
        (the ANDES backend is single-threaded). The full
        seed-sweep + ranker re-evaluation reported here cost
        $\approx 14$\,h. Resource requirements are modest by ML
        standards but non-trivial for a research lab without GPU.
  \item \textbf{Code-drift caveat.} The recurrent agent suffers a
        $\sim$19\% headline regression between two reproductions of
        the same training configuration nine rounds apart (CLM-0104):
        seed 51 LSTM at $\tau = 0.005$, warm-up 5 scored
        $v_{6\text{-axis}} = 0.526$ in our R57-$\alpha$ baseline and
        $0.426$ when re-trained nine rounds later under what we
        believed was the same code path. The memoryless TD3 and
        SAC variants reproduce cleanly under the same nine-round
        gap. The candidate commit range was bisected to six
        commits between rounds R58 and R65; the leading suspect is
        the R61 monitor-extension commit, which added a Python
        atexit handler that may interact with PyTorch's CUDA-less
        RNG. We did not complete the bisection in time for this
        submission. The numbers in this paper are pinned to the
        post-regression code path. For reproducibility we publish
        (a) the git SHA of the commit producing each table (see
        \texttt{README.md}) and (b) the bit-identical regression
        test \texttt{tests/test\_v4\_env\_regression.py} that has
        held the V4 environment to $10^{-9}$ tolerance against the
        pre-refactor baseline. Whether the LSTM-only $\sim$19\%
        drift is a fundamental reproducibility limit or an
        unidentified bug remains open.
\end{enumerate}

% ─────────────────────────────────────────────────────────────────────
\section{Conclusion}\label{sec:conclusion}

We have presented three contributions to the multi-agent VSG
inertia/damping control problem on the modified Kundur benchmark:
an eleven-axis multiplicative-gating ranker that catches the
single-axis failure modes invisible to the prior six-axis geometric
mean; a recurrent TD3+LSTM multi-agent policy with an episode-keyed
sequence replay that achieves a new single-seed SOTA of v3.1 =
0.430; and a clean negative result for inference-time HAWE
ensembling on the recurrent policy class, attributable to
hidden-state divergence.

The headline numbers should be read as: the policy class matters
(LSTM gains $+9\%$ over the strongest feed-forward variant), the
evaluation metric matters (the v3.1 ranking differs from v3.0 by
2--3 places on the leaderboard, Table~\ref{tab:demotion}), and the
seed filter matters (the unfiltered 12-seed mean is roughly
$0.5\times$ the healthy 6-seed mean because the five drift seeds
contribute v$_{3.1} \leq 0.07$ each; the magnitude depends on the
seed pool but the direction is robust).

Future work includes: (i) a hidden-state-mixing HAWE variant to
test whether the recurrent ensemble regression is fundamental or
protocol-specific; (ii) extension of the eleven-axis ranker
boundary to a learned weighting rather than a uniform geometric
mean; (iii) replication of the LSTM result on the modified NE-39
benchmark from \cite[Sec.~IV-G]{yang2023ddic}, conditional on
resolving the $M_0 < 20$ TDS-divergence that prevented our prior
NE-39 attempts; (iv) verification that the policy-class ranking
\emph{persists at the paper's full 500-episode convergence
horizon} -- our reported per-seed runs are 75 episodes (the
empirical stabilisation point on the ANDES backend; see
Sec.~\ref{sec:setup}), and the paper used 2000 (open question
Q-0008 in the project ledger); and (v) isolation of the LSTM-only
$\sim$19\% code-drift root cause (Sec.~\ref{sec:limits}).

\section*{Reproducibility}

All training scripts, evaluation traces, and the v3.1 ranker source
are released. The complete claim ledger covering rounds R67--R77 is
under \texttt{memory/claims/CLM-\{0086,...,0136\}.md} with round
verdicts under \texttt{memory/rounds/R\{67--77\}/verdict.md}. The
SHA of the commit producing this paper's tables is recorded in the
paper's \texttt{README.md}.

% ─────────────────────────────────────────────────────────────────────
\bibliographystyle{IEEEtran}
\bibliography{refs}

\end{document}
